\chapter{Running Simulations}
\label{ch:simulations}

\section{The calculation dialog}

\menu{Simulation}\menusep\menu{New Calculation (ASE)} (\kbd{Ctrl}+\kbd{R})
is the main entry point. The left half is a form; the right half is the ASE
script that form produces, live and editable.

\subsection{Calculators}

\begin{center}
\small
\begin{tabular}{@{}lp{8.4cm}@{}}
\toprule
Calculator & Notes \\
\midrule
EMT             & Effective medium theory. Fast, ships with ASE, useful for
                  testing a workflow before committing compute. \\
Lennard-Jones   & Ships with ASE. \\
Quantum ESPRESSO& DFT; needs \texttt{pw.x} and pseudopotentials. \\
VASP            & DFT; needs a license and POTCARs. \\
MACE            & Machine-learning potential; needs \texttt{mace-torch}. \\
GPAW            & DFT as a Python package. \\
SIESTA          & DFT; needs the \texttt{siesta} binary and pseudopotentials. \\
ORCA            & Quantum chemistry; needs the ORCA binaries. \\
\bottomrule
\end{tabular}
\end{center}

The DFT entries generate working script skeletons with clearly marked
\texttt{EDIT ME} lines where you must supply pseudopotential paths or
launch commands.

\subsection{Tasks}

\begin{description}
  \item[\ui{Single-point energy}] One energy and force evaluation.
  \item[\ui{Geometry optimization (BFGS)}] Controlled by
    \ui{Force convergence (fmax)} (default 0.05\,eV/\AA) and
    \ui{Max optimization steps} (default 200). Writes \file{opt.traj}.
  \item[\ui{Molecular dynamics}] See below. Writes \file{md.traj}.
\end{description}

\subsection{Molecular dynamics ensembles}

\begin{center}
\small
\begin{tabular}{@{}llp{5.6cm}@{}}
\toprule
Ensemble & Thermostat / barostat & Key parameters \\
\midrule
NVE  & Velocity Verlet                    & timestep \\
NVT  & Langevin (default)                 & temperature, friction \\
NVT  & Andersen                           & temperature, collision probability \\
NVT  & Berendsen                          & temperature, $\tau_T$ \\
NVT  & Nos\'e--Hoover chain               & temperature, $\tau_T$ \\
NPT  & Berendsen                          & temperature, pressure, $\tau_T$, $\tau_P$ \\
NPT  & Nos\'e--Hoover / Parrinello--Rahman & temperature, pressure, $\tau_T$, $\tau_P$ \\
\bottomrule
\end{tabular}
\end{center}

Shared parameters: \ui{Temperature} (default 300\,K), \ui{Timestep}
(default 1\,fs), \ui{MD steps} (default 1000),
\ui{Thermostat coupling} (default 100\,fs), \ui{Barostat coupling}
(default 1000\,fs) and \ui{Pressure} (default 0\,GPa). Fields enable
themselves only for the ensembles that use them.

\begin{calnote}
Every MD run initialises velocities from a Maxwell-Boltzmann distribution
and then removes the net centre-of-mass momentum, so the system does not
drift as a whole. For isolated (non-periodic) systems the net angular
momentum is removed as well. Without this, a slow overall translation
contaminates every subsequent analysis --- diffusion coefficients most of
all.
\end{calnote}

\subsection{DFT and ORCA parameters}

DFT calculators expose \ui{Plane-wave cutoff} (default 550\,eV) and a
\ui{k-point grid} (default $4\times4\times4$).

ORCA has its own group:

\begin{description}
  \item[\ui{ORCA method}] Editable: \texttt{B3LYP}, \texttt{PBE0},
    \texttt{r2SCAN}, \texttt{wB97X-D3}, \texttt{HF}, \texttt{MP2}, or any
    ORCA keyword you type.
  \item[\ui{ORCA basis}] Editable: \texttt{def2-SVP}, \texttt{def2-TZVP},
    \texttt{def2-TZVPP}, \texttt{cc-pVDZ}, \texttt{cc-pVTZ}, \ldots
  \item[\ui{Charge}] $-10$ to 10.
  \item[\ui{Multiplicity}] $2S+1$, 1 to 11.
  \item[\ui{Solvation}] \ui{None (gas phase)}, \ui{CPCM} or \ui{SMD}, with
    \ui{Solvent} naming the medium (water, acetonitrile, toluene, \ldots).
\end{description}

ASE writes the ORCA \file{.inp} file into the job directory and parses the
output. The generated script contains an \texttt{EDIT ME} line for the path
to your \texttt{orca} binary.

\subsection{MACE models}

\begin{description}
  \item[\ui{MACE model}] \ui{MACE-MP-0 (universal, materials)},
    \ui{MACE-OFF (universal, organic molecules)}, or
    \ui{Custom trained model (file)}.
  \item[\ui{MACE model size}] \texttt{small}, \texttt{medium} (default),
    \texttt{large}.
  \item[\ui{Custom model file}] A \file{.model} or \file{.pt} checkpoint.
  \item[\ui{MACE device}] \texttt{cpu}, \texttt{cuda} or \texttt{mps}.
\end{description}

Foundation models download automatically on first use and are cached in
\file{\textasciitilde/.cache/mace}.

\section{The script editor}

The right-hand pane always shows the complete, standalone ASE script the
form describes --- syntax-highlighted and fully editable.

Start typing in it and Calango shows \emph{edited --- form sync paused} in
amber: from then on the form no longer overwrites your text, so hand edits
are never silently lost. \ui{Regenerate} discards them and re-syncs.
\ui{Save Script\ldots} writes the current text to a \file{.py} file.

\begin{caltip}
This is the intended escape hatch for anything the GUI does not expose.
Configure what you can in the form, then edit the script for the rest ---
constraints, custom observers, a different optimiser. The script is
ordinary ASE and runs unmodified on a cluster.
\end{caltip}

\section{Execution environments}
\label{sec:exec-env}

The \ui{Execution Environment} group decides which Python actually runs the
job --- independent of the interpreter Calango embeds.

\begin{description}
  \item[Leave it empty] The job uses the embedded interpreter. The status
    line names it.
  \item[\ui{Env Folder\ldots}] Point at a Conda environment directory;
    Calango finds \file{bin/python} inside it.
  \item[\ui{Python\ldots}] Point directly at an interpreter.
\end{description}

The status line confirms the resolution live, turning red when no
interpreter can be found at the given path. The setting is remembered
between sessions, and the same selector appears in the phonon builder and
the band-structure dialog.

\begin{caltip}
This is how you keep heavyweight solver stacks isolated: a GPAW environment
for band structures, a \texttt{mace-torch} environment with CUDA for ML
potentials, and a minimal ASE environment for Calango itself.
\end{caltip}

\section{The Job panel}
\label{sec:job-panel}

Jobs run as separate processes in their own directory, so a crashing solver
cannot take the GUI with it. The \ui{Job} dock (zone 10) has five tabs.

\subsection{Log}

A status label (\emph{Idle}, \emph{Running\ldots}, \emph{Finished (exit
N)}, \emph{Crashed}), a progress bar, a \ui{Kill} button, and the live
output. Standard error is coloured red, the start line blue, and a clean
exit green.

\subsection{Metric plots}

Four tabs plot job observables live, parsed from markers the generated
scripts print:

\begin{center}
\small
\begin{tabular}{@{}llll@{}}
\toprule
Tab & Quantity & Unit & Reference line \\
\midrule
\ui{Energy}      & Total energy    & eV      & --- \\
\ui{Temperature} & Ionic temperature & K     & thermostat setpoint \\
\ui{Force}       & Max $|F|$       & eV/\AA  & --- \\
\ui{Pressure}    & Scalar pressure & GPa     & barostat setpoint \\
\bottomrule
\end{tabular}
\end{center}

Each plot shows the running series with the latest value read out
numerically, and a dashed reference line at the setpoint for thermostatted
or barostatted runs --- so drift away from the target is visible at a
glance. Every tab has its own \ui{Export Data\ldots} button writing CSV or
whitespace-separated \file{.dat}. The series are saved in the project file,
so a reopened project shows the plots from the run that produced it.

\section{Live trajectory streaming}
\label{sec:live-streaming}

MD runs and geometry relaxations do not make you wait. The moment the job
starts, Calango opens a trajectory tab marked \emph{(live)} and streams
each newly computed geometry into it as the simulation produces it.

\begin{itemize}
  \item The timeline grows as frames arrive, and follows the newest frame
    automatically.
  \item Scrub back to an earlier frame and the view stays there --- Calango
    stops following so you can inspect while the run continues.
  \item Everything else works normally on the partial trajectory:
    measurement, colour mapping, representation changes.
  \item When the job finishes the \emph{(live)} marker is dropped and the
    tab becomes an ordinary trajectory.
\end{itemize}

Relaxations stream every step; MD streams at an interval chosen so a run
produces at most about 400 streamed frames, keeping long simulations
responsive. The complete trajectory is still written to
\file{opt.traj} or \file{md.traj} in the job directory.

\section{What happens when a job finishes}

Calango inspects the job directory and does the useful thing:

\begin{itemize}
  \item A live-streamed run keeps the tab it already built.
  \item An electronic-structure run opens the band/PDOS viewer
    (Chapter~\ref{ch:electronic}).
  \item A trajectory (\file{md.traj}, \file{opt.traj}) opens in a new tab.
  \item Otherwise, if a final geometry exists
    (\file{optimized.extxyz}), Calango offers to load it.
\end{itemize}

The process manager entry switches to \emph{completed} or \emph{failed},
and its directory link remains available for the rest of the session and
beyond.

\section{The MLIP dataset manager}
\label{sec:dataset-manager}

\menu{Simulation}\menusep\menu{Dataset Manager (MLIP)} assembles training
datasets for machine-learning interatomic potentials from the trajectories
you have generated --- MD runs, noise ensembles, relaxation paths --- and
partitions them reproducibly.

\subsection{Loading structures}

\ui{Add Files\ldots} accepts multiple files of mixed provenance:
\file{.extxyz}, \file{.traj}, \file{.cif}, VASP files, ASE databases. Each
is listed with its frame count and the chemical formulas it contains, so a
heterogeneous dataset stays legible. The summary line reports the total.

\subsection{Splitting}

\begin{description}
  \item[\ui{Training}, \ui{Validation}] Percentages; \ui{Test} is the
    remainder, shown live.
  \item[\ui{Random seed}] Default 42.
\end{description}

The split is deterministic: the same frame count, percentages and seed
always produce exactly the same partition, and the three subsets are
disjoint and cover every frame.

\subsection{Query-by-Committee ensembles}

Active learning needs an ensemble of models trained on different data, so
their disagreement can be used as an uncertainty estimate.

\begin{description}
  \item[\ui{Committee members}] 1 exports a single dataset; $N > 1$ adds
    \file{committee\_01}\ldots\file{committee\_N} subdirectories, each with
    its own training set.
  \item[\ui{Diversity}] How the members are made to differ:
    \begin{itemize}
      \item \ui{Independent splits (seed + k)} re-splits the non-test pool
        with a different seed per member, keeping one common test set.
      \item \ui{Bootstrap resampling of the training set} draws each
        member's training set with replacement --- the classical bagging
        construction, which reuses roughly 63\,\% of the original frames
        per member.
    \end{itemize}
\end{description}

\subsection{Export}

\begin{description}
  \item[\ui{Extended XYZ (MACE-ready)}] Writes \file{train.extxyz},
    \file{valid.extxyz} and \file{test.extxyz} (plus committee
    subdirectories), the layout MACE and most training scripts expect.
  \item[\ui{ASE database (.db)}] One SQLite database per subset.
\end{description}

\begin{calnote}
Export goes through \texttt{ase.io} directly rather than Calango's internal
structure model, specifically so that energies, forces and stresses
attached to the frames survive into the training files. A dataset exported
from a finished MD run carries the labels a potential needs to train on.
\end{calnote}
